% !PRINTED 315
\weberpage{315}
%<<CONT
so sind $A_1^2, A_2, A_3^2, A_4, A_5^2$ nach den Grundsätzen des §\,73
\gesperrt{ganze rationale Funktionen} von

\weq{24}{u^{24} + \frac{2^{12}}{u^{24}}.}

Die Entwickelung von (24) nach steigenden Potenzen von $q$ fängt an mit
$q^{-1}$, während der Anfang der Entwickelung von $w_z$

\[
q^{-\frac{1}{10}}e^{-\frac{c\pi i}{60}}
\]

ist. Die Größen $A_1^2, A_2, A_3^3, A_4$ können % [sic] Gedruckt ist hier A_3^3; oben steht A_3^2.
 daher nicht einmal die
erste Potenz von (24) enthalten und sind also konstant, während $A_5^2$
die Form hat:

\[
A_5^2 = u^{24} + \frac{2^{12}}{u^{24}} + C,
\]

worin $C$ eine Konstante ist. Die Werte der Konstanten kann man
bestimmen, wenn man die Werte der $w_z$ für $\omega = i$ kennt. Es ist
aber für $\omega = i$ [vgl. oben (6), (7)]:

\[
\begin{aligned}
u &= f(i) = \sqrt[4]{2},\\
v_3 &= f\left(\frac{i+48}{5}\right) = f\left(\frac{i-2}{5}+10\right) = e^{-\frac{10\pi i}{24}}f\left(\frac{i-2}{5}\right),\\
&= e^{-\frac{10\pi i}{24}}f(i+2) = -i\sqrt[4]{2},\\
v_2 &= i\sqrt[4]{2}
\end{aligned}
\]

und folglich nach (18):

\[
\begin{aligned}
v_\infty &= v_0 = \sqrt[4]{2}\,\frac{1+\sqrt{5}}{2},\\
v_1 &= v_4 = \sqrt[4]{2}\,\frac{1+\sqrt{5}}{2}. % [sic] Beide gedruckten Zeilen haben dasselbe Pluszeichen.
\end{aligned}
\]

Folglich wird:

\[
w_0 = 0, \quad w_1 = w_2 = i\sqrt{5}, \quad w_3 = w_4 = -i\sqrt{5}.
\]

Danach wird zunächst $C = -2^7$, also

\[
A_5^2 = \left(u^{12} - \frac{64}{u^{12}}\right)^2,
\]

und daraus durch Vergleichung der Anfänge der Entwickelung

\[
A_5 = -u^{12} + \frac{64}{u^{12}},
\]

und die Gleichung für $w$ bekommt die Form:

\weq{25}{w(w^2+5)^2 = u^{12} - \frac{64}{u^{12}}.}
